% Lösungen Einheit 2
\einheitenkopf{Einheit 2 von 4 · Satz vom Nullprodukt}
\zweigzeile{Hier lernst du, Gleichungen wie $(x - a)\cdot(x - b) = 0$ mit dem Satz vom Nullprodukt zu lösen, auch durch Ausklammern von $x$ · neu oder schon bekannt – je nach Buch (an der Oberschule erst in Klasse 10) · P10 · baut auf: lineare Gleichungen, Ausklammern, Einsetzen prüfen}
\erg{19}{a) gilt \quad b) gilt nicht \quad c) gilt \quad d) gilt nicht \quad e) gilt}
\erg{20}{a) $x_1 = 1$, $x_2 = 4$ \quad b) $x_1 = 3$, $x_2 = 6$ \quad c) $x_1 = 8$, $x_2 = 2$ \quad d) $x_1 = 5$, $x_2 = 7$ \quad e) $x_1 = -4$, $x_2 = 1$ \quad f) $x_1 = 0$, $x_2 = -6$ \quad g) eine Lösung: $x = -3$ \quad h) $(-5)\cdot(-1) = 5$, $5 = 0$ (fA)}
\erg{21}{a) $x\cdot(x + 4) = 0$; $x_1 = 0$, $x_2 = -4$ \quad b) $x_1 = 0$, $x_2 = -9$ \quad c) $x_1 = 0$, $x_2 = 8$ \quad d) $2x\cdot(x + 5) = 0$; $x_1 = 0$, $x_2 = -5$ \quad e) $x^2 + x - 30 = 0$ \quad f) $2\cdot(x^2 - x - 12) = 0$, $2\cdot(x - 4)\cdot(x + 3) = 0$; $x_1 = 4$, $x_2 = -3$ \quad g) $-3$, denn $(-3)\cdot 4 = -12$}
\erg{22}{a) Fehler: durch $x$ geteilt – $x$ kann null sein, die Lösung $x = 0$ geht verloren; richtig: $x^2 - 6x = 0$, $x\cdot(x - 6) = 0$, $x_1 = 0$, $x_2 = 6$ \quad b) $4x^2 - 28x = 0$, $4x\cdot(x - 7) = 0$; $x_1 = 0$, $x_2 = 7$}
\erg{23}{a) für $x = 4$ ist der erste Faktor null, für $x = -1$ der zweite – ein Produkt ist null, sobald ein Faktor null ist \quad b) nein: der Satz gilt nur, wenn auf einer Seite null steht; Probe $x = 6$: $6\cdot 5 = 30$, nicht $6$; richtig: $x^2 - x - 6 = 0$, $x_1 = 3$, $x_2 = -2$ (Formel aus Einheit 3)}
\erg{24}{a) $x\cdot(x - 48) = 0$; $x_1 = 0$, $x_2 = 48$; Spannweite $48\,\text{m}$ \quad b) Mitte bei $x = 24$: $y = -0{,}05\cdot 24\cdot(-24) = 28{,}8$; der Bogen ist dort $28{,}8\,\text{m}$ hoch}
